%% kapa: exam=1 round=36 year=2025 period=1 form=A subject=civil_law subject_ko=민법 seq=38 number=38
%% source: Q-Net 감정평가사 기출문제 (cst00309) artlSeq=5231525 file=raw/1008/5231525/2025년 제36회 감정평가사 1차시험 1교시 문제지 원본.pdf
%% conversion: kordoc 4.24.51 -> tools/gen1.py (auto); review=proofread-sonnet; flags=none
\begin{question}{38}
채권질권에 관한 설명으로 옳지 않은 것은? (다툼이 있으면 판례에 따름)
\begin{choices}
  \item 임차보증금반환청구권은 채권질권의 목적이 될 수 있다.
  \item 채권증서가 있는 지명채권의 경우 그 증서를 질권자에게 교부해야 질권의 효력이 생긴다.
  \item 채권의 목적물이 금전 이외의 물건인 때에는 채권질권자는 직접 변제받은 물건에 대해 질권을 행사할 수 있다.
  \item 질권설정자가 질권이 설정된 채권을 양도하는 경우 질권자의 동의는 필요 없다.
  \item 질권설정자와 제3채무자가 질권자의 동의 없이 질권이 설정된 채권을 소멸하게 하는 행위를 한 경우 특별한 사정이 없는 한 제3자도 그 무효를 주장할 수 있다.
\end{choices}
\end{question}
\answer{5}
